\begin{proposition}[Proved: tame ramification above positive-genus exceptional curves]
\label{D:lc-elliptic-tame}
Let $(S,s)$ be a normal log canonical surface singularity over an
algebraically closed field $k$.  Let $E$ be an exceptional divisorial
valuation over $s$ whose residue curve has positive geometric genus.
Every finite separable quasi-\'etale cover $S'\to S$ is tame at $E$:
for every extension $F$ of this valuation, its ramification index is
prime to the characteristic exponent of $k$ and
$\kappa(F)/\kappa(E)$ is separable.
In the application below we verify directly the log canonical and
Gorenstein properties of the contraction of one smooth elliptic curve,
so no classification in small characteristic is required.
\end{proposition}
\begin{proof}
We use surface resolution and point-blowup factorization in arbitrary
characteristic, as stated in the Stacks Project,
Theorem~54.14.5 and Lemma~54.17.1, tags 0BGP and 0C5R.
Contraction of a smooth rational exceptional curve of self-intersection
$-1$ is Lemma~54.16.8, tag 0C2M
\cite{StacksDSurfaceResolution}.
The exceptional intersection matrix on a resolution over an affine
neighborhood of the singularity is negative definite by
\cite[section 10.1, Theorem 10.1]{Kollar2013D}; the source explicitly
works with excellent surfaces in arbitrary characteristic in that
section.

We first give the ramification calculation, including the case of an
inseparable residue extension.  Let $F$ be a divisorial valuation in
a finite separable extension of the surface function field, and let
$E$ be its restriction.  Choose models regular at the respective
generic points and a morphism between neighborhoods of those points.
These neighborhoods are smooth over the perfect field $k$.
Choose a uniformizer $t$ at $E$ and a lift $z$ of a separating
transcendence parameter of $\kappa(E)/k$; then $dt\wedge dz$
generates the canonical module at $E$.  Write
$t=u s^e$ at $F$, with $u$ a unit and $s$ a uniformizer.
Choose similarly a separating residue parameter $w$ at $F$.
The order $\delta_F$ of the determinant of differentials is the
coefficient of the relative canonical divisor, and
\[
 d(u s^e)\wedge dz
   =s^{e-1}\bigl(eu\,ds\wedge dz+s\,du\wedge dz\bigr).
\]
Consequently $\delta_F\geq e-1$.  Equality holds precisely when
$e$ is nonzero in $k$ and $d\bar z$ is nonzero in
$\Omega^1_{\kappa(F)/k}$.  For a finite extension of one-variable
function fields over a perfect field, the latter condition is
equivalent to separability of $\kappa(F)/\kappa(E)$.
Thus equality is precisely tame ramification, including the
residue-separability requirement.  The determinant is not identically
zero because the extension of surface function fields is separable.

Since $S'\to S$ is quasi-\'etale, its canonical divisor satisfies
$K_{S'}=\pi^*K_S$ as a rational Cartier divisor: the equality holds
at every codimension-one point, and normality extends the equality
of Weil divisors.  In particular $S'$ is $\mathbb Q$-Gorenstein.
On models extracting a pair of valuations as above, comparison of
canonical divisors gives the exact discrepancy formula
\[
 a(F,S')=e\,a(E,S)+\delta_F.
\]
The restriction of a divisorial valuation through a finite field
extension is again divisorial.  Applying the lower bound to every
such valuation shows that $S'$ is log canonical wherever $S$ is:
$a(F,S')\geq-e+(e-1)=-1$.

We next prove the needed fact about positive-genus exceptional curves
without a classification theorem.  On a minimal resolution of an
arbitrary log canonical surface singularity, write
\[
 K_Z=f^*K_S+\sum_i a_i E_i,\qquad M=(E_i\cdot E_j).
\]
For each exceptional integral curve, adjunction gives
$K_Z\cdot E_i=2p_a(E_i)-2-E_i^2\geq0$.
Indeed $E_i^2<0$; when $p_a(E_i)=0$, the curve is smooth rational
and minimality excludes $E_i^2=-1$.  Thus $Ma=b$ with $b\geq0$.
All off-diagonal entries of $M$ are nonnegative.  If the positive
part $a^+$ were nonzero, negative definiteness and the disjoint
supports of $a^+$ and the nonpositive part $a^-$ would imply
\[
 0\leq (a^+)^{\mathsf T}b
    =(a^+)^{\mathsf T}Ma^++(a^+)^{\mathsf T}Ma^-<0.
\]
Hence every $a_i\leq0$.  If the normalization of $E_i$ has positive
genus, then $p_a(E_i)\geq1$, so
\[
 -E_i^2\leq K_Z\cdot E_i
      =\sum_j a_j(E_j\cdot E_i)
      \leq a_iE_i^2\leq-E_i^2.
\]
The last inequality uses log canonicity $a_i\geq-1$.
Equality forces $p_a(E_i)=1$ and $a_i=-1$.
A positive-genus valuation appearing on any other resolution
survives the contractions to a minimal resolution, because the
curves contracted between smooth surfaces are rational.  Its
discrepancy is therefore also $-1$.

Return to the specified $E$ and a valuation $F$ above it.
The residue curves admit a finite dominant map after normalization.
The curve of $F$ cannot be rational, since a nonconstant map from
$\mathbb P^1$ to a positive-genus smooth proper curve is impossible,
also for inseparable maps.  Equivalently we may apply L\"uroth's
theorem to the corresponding function fields.  Both $E$ and $F$
therefore have discrepancy $-1$ on the respective log canonical
surfaces.  The discrepancy formula now gives
$-1=-e+\delta_F$, hence $\delta_F=e-1$.
The differential calculation proves the asserted tameness.
\end{proof}

\begin{corollary}[Proved: the elliptic three-point surface with trivial Picard group]
\label{D:schroer-elliptic}
In Schr\"oer's construction of a normal proper surface $S$ with
$\operatorname{Pic}(S)=0$, take the initial curve $C$ to be elliptic
over an algebraically closed field.  Let $s_0,s_1,s_2$ be its three
contracted points.  Then
$\pi_1(S\setminus\{s_0,s_1,s_2\})$ is topologically finitely
generated, in every characteristic.
\end{corollary}
\begin{proof}
We use exactly the construction in
\cite[sections 2.5 and 3.1--3.2, author's PDF pp.~4--5]{Schroer1999D}.
Choose $c_1,c_2\in C$ as in that source, blow up
$\mathbb P^1\times C$ at $(b_0,c_1),(b_1,c_1),(b_2,c_2)$,
and call the resulting smooth projective surface $X$.
Write $R_i$ for the strict vertical fibres, $A$ for the strict
horizontal curve through $c_1$, and $B$ for that through $c_2$.
The disjoint connected configurations $R_0$, $R_1$ and $A\cup R_2$
contract to $s_0,s_1,s_2$.  The common open in question is
\[
 U=X\setminus(R_0\cup R_1\cup A\cup R_2).
\]

We check the local hypothesis at $s_0,s_1$ explicitly.  For
$i\in\{0,1\}$ choose a coordinate $t$ on $\mathbb P^1$ vanishing
at $b_i$, and a nowhere-vanishing regular differential $\eta$ on
the elliptic curve $C$.  The rational form
$(dt/t)\wedge\eta$ has a simple pole on the fibre through $b_i$.
Blowing up $(b_i,c_1)$ adds discrepancy $1$ and pulls back that
pole with multiplicity $1$, so the exceptional rational curve has
coefficient zero.  On a neighborhood of $R_i$ the divisor of the
pulled-back form is exactly $-R_i$.  Its other divisor components
have closed images avoiding $s_i$, so we can shrink around $s_i$.
The descended form has divisor zero on the punctured neighborhood;
normality identifies its reflexive canonical sheaf with
$\mathcal O_S$ there and across $s_i$.  Thus $K_S$ is Cartier near
$s_i$, and the displayed resolution has the single discrepancy
$-1$ on the smooth elliptic curve $R_i$.  It is log canonical:
the boundary consists of that one smooth curve, and subsequent
point blowups preserve discrepancies at least $-1$.

Section 3.1 of the same source constructs a \emph{projective}
normal surface $T$ by contracting only $A\cup R_2$.  Explicitly,
the effective divisor $D=3B+2R_0+2R_1$ satisfies
$D\cdot B=D\cdot R_0=D\cdot R_1=1$, and the associated contraction
is the one asserted there.  Hence
$U=T\setminus(R_0\cup R_1\cup\{q\})$, where $q$ is the image of
$A\cup R_2$.

Every connected finite \'etale cover of $U$ extends by normalization
to a finite cover of $S$ that is \'etale away from its three
omitted points, hence quasi-\'etale.  Proposition~\ref{D:lc-elliptic-tame}
shows that this cover is tame at the two divisorial boundary
components $R_0,R_1$ on $T$.  The remaining boundary point $q$
has codimension two and imposes no condition in $Q(T,T\setminus U)$.
Thus every finite \'etale cover of $U$ already belongs to this
fixed-boundary tame category, and
\[
 \pi_1(U)=Q(T,T\setminus U).
\]
The right-hand side is finitely generated by
Proposition~\ref{D:fixed-boundary-positive}, since $T$ is projective.
We do not require the singularity at $q$ to be log canonical.
Indeed, in the elliptic case the discrepancies on the configuration
$A\cup R_2$ are $-1$ and $-2$, respectively; leaving it as a
codimension-two boundary point is essential to this proof.
\end{proof}
